\documentclass{article}
% The preamble every project's paper shares: % The preamble every project's paper shares: % The preamble every project's paper shares: \input{../../preamble} from
% projects/<project>/paper/. Boards come from chessboard and xskak, plots from
% pgfplots; every number a paper prints comes from its generated.tex.
\usepackage[LSBC4,T1]{fontenc}
\usepackage[utf8]{inputenc}
\usepackage{lmodern}
\usepackage[margin=1in]{geometry}
\usepackage{microtype}
\usepackage{amsmath}
\usepackage{booktabs}
\usepackage{tabularx}
\usepackage{array}
\usepackage{xcolor}
\usepackage[obeyspaces,spaces]{url}
\usepackage{xskak}
\usepackage{pgfplots}
\pgfplotsset{compat=1.18}
\usepackage{caption}
\usepackage{subcaption}
\usepackage{enumitem}
\usepackage[section]{placeins}
\usepackage[hidelinks]{hyperref}

\setchessboard{boardfontencoding=LSBC4, boardfontfamily=skaknew, blackfieldcolor=black!22,
  whitefieldcolor=white, setfontcolors, showmover=false, label=true, labelfontsize=6pt,
  boardfontsize=14pt}

% A move or line in UCI, and a position in FEN, both breakable at a slash.
\newcommand{\uci}[1]{\texttt{#1}}
\newcommand{\fen}[1]{{\footnotesize\path{#1}}}
\newcolumntype{L}{>{\raggedright\arraybackslash}X}

\setlength{\parskip}{0.4em}
\setlength{\parindent}{0pt}
 from
% projects/<project>/paper/. Boards come from chessboard and xskak, plots from
% pgfplots; every number a paper prints comes from its generated.tex.
\usepackage[LSBC4,T1]{fontenc}
\usepackage[utf8]{inputenc}
\usepackage{lmodern}
\usepackage[margin=1in]{geometry}
\usepackage{microtype}
\usepackage{amsmath}
\usepackage{booktabs}
\usepackage{tabularx}
\usepackage{array}
\usepackage{xcolor}
\usepackage[obeyspaces,spaces]{url}
\usepackage{xskak}
\usepackage{pgfplots}
\pgfplotsset{compat=1.18}
\usepackage{caption}
\usepackage{subcaption}
\usepackage{enumitem}
\usepackage[section]{placeins}
\usepackage[hidelinks]{hyperref}

\setchessboard{boardfontencoding=LSBC4, boardfontfamily=skaknew, blackfieldcolor=black!22,
  whitefieldcolor=white, setfontcolors, showmover=false, label=true, labelfontsize=6pt,
  boardfontsize=14pt}

% A move or line in UCI, and a position in FEN, both breakable at a slash.
\newcommand{\uci}[1]{\texttt{#1}}
\newcommand{\fen}[1]{{\footnotesize\path{#1}}}
\newcolumntype{L}{>{\raggedright\arraybackslash}X}

\setlength{\parskip}{0.4em}
\setlength{\parindent}{0pt}
 from
% projects/<project>/paper/. Boards come from chessboard and xskak, plots from
% pgfplots; every number a paper prints comes from its generated.tex.
\usepackage[LSBC4,T1]{fontenc}
\usepackage[utf8]{inputenc}
\usepackage{lmodern}
\usepackage[margin=1in]{geometry}
\usepackage{microtype}
\usepackage{amsmath}
\usepackage{booktabs}
\usepackage{tabularx}
\usepackage{array}
\usepackage{xcolor}
\usepackage[obeyspaces,spaces]{url}
\usepackage{xskak}
\usepackage{pgfplots}
\pgfplotsset{compat=1.18}
\usepackage{caption}
\usepackage{subcaption}
\usepackage{enumitem}
\usepackage[section]{placeins}
\usepackage[hidelinks]{hyperref}

\setchessboard{boardfontencoding=LSBC4, boardfontfamily=skaknew, blackfieldcolor=black!22,
  whitefieldcolor=white, setfontcolors, showmover=false, label=true, labelfontsize=6pt,
  boardfontsize=14pt}

% A move or line in UCI, and a position in FEN, both breakable at a slash.
\newcommand{\uci}[1]{\texttt{#1}}
\newcommand{\fen}[1]{{\footnotesize\path{#1}}}
\newcolumntype{L}{>{\raggedright\arraybackslash}X}

\setlength{\parskip}{0.4em}
\setlength{\parindent}{0pt}

% Written by projects/maia-player-offsets/paper/generate.mjs; do not edit by hand.

\newcommand{\Phase}{1}
\newcommand{\ModelName}{Maia-2}
\newcommand{\AppCommit}{\texttt{f4bd63f19b343f686fb97ea4cf70b211c0bddc0a}}
\newcommand{\RecordsCommit}{\texttt{aed714ada579782ce73630596e59c8b65e74cb30}}
\newcommand{\Games}{648}
\newcommand{\Moves}{33522}
\newcommand{\FitGames}{324}
\newcommand{\FitMoves}{17055}
\newcommand{\HeldGames}{324}
\newcommand{\HeldMoves}{16467}
\newcommand{\EmptyGames}{8}
\newcommand{\FrozenGames}{60}
\newcommand{\GuardGames}{28}
\newcommand{\RecordedEvents}{328}
\newcommand{\RecordedExact}{244}
\newcommand{\Base}{250}
\newcommand{\ACount}{441}
\newcommand{\BCount}{147}
\newcommand{\OffsetMin}{0}
\newcommand{\OffsetMax}{500}
\newcommand{\OffsetStep}{25}
\newcommand{\Anchor}{1000}
\newcommand{\SlopeList}{$-$0.75, $-$0.50, $-$0.25, 0.00, 0.25, 0.50, 0.75}
\newcommand{\Draws}{1000}
\newcommand{\Seed}{7}
\newcommand{\GuardMargin}{0.005}
\newcommand{\OrdinaryCI}{95}
\newcommand{\AdjustedCI}{97.5}
\newcommand{\Alpha}{0.05}
\newcommand{\GuardCells}{4}
\newcommand{\MinPly}{11}
\newcommand{\MinClock}{30}
\newcommand{\ErrorLoss}{5}
\newcommand{\PolicyCalls}{900016}
\newcommand{\SampleMoves}{32}
\newcommand{\SampleCalls}{890}
\newcommand{\PolicySeconds}{10.27}
\newcommand{\CallMS}{21.03}
\newcommand{\Startup}{0.374}
\newcommand{\ReviewGames}{2}
\newcommand{\ReviewPlies}{140}
\newcommand{\ReviewSeconds}{108.40}
\newcommand{\FullHours}{2.99}
\newcommand{\PlanningHours}{4.49}
\newcommand{\Processes}{2}
\newcommand{\Nice}{10}
\newcommand{\Allowance}{1.5}
\newcommand{\StorageMB}{841}
\newcommand{\WeightsMB}{84.95}
\newcommand{\CostTable}{%
\begin{tabular}{lrr}\toprule Step & Work & Projected hours \\ \midrule
Move policies & 900016 & 2.63 \\
Routine reviews & 28 games & 0.36 \\
Punish policies (old-event estimate) & 1927 & 0.006 \\
\bottomrule\end{tabular}
}

\title{Does Maia need a different offset for the player?\\
\large A side split and a rating function on held-out games}
\author{Chess Coach research lab}
\date{}
\begin{document}
\maketitle

\begin{abstract}
The app adds \Base{} points to a chess.com rating before choosing a
\ModelName{} band. Its predecessor kept that offset, but found exploratory
differences by side and rating. We compare separate offsets for the player
and his opponents with a single offset function of rating. The fits use
\FitMoves{} moves; the test uses \HeldMoves{} moves in \HeldGames{} games.
Every held-out game also receives a routine review for the punish guard.
\Verdict{}
\end{abstract}

\section{The data and the question}
We measure app commit \AppCommit. The research records are at
\RecordsCommit. The archive contributes \Moves{} moves in \Games{} games.
The predecessor's position script is run with only its private output
directory changed. Moves count from ply \MinPly{} when both clocks at the
position have at least \MinClock{} seconds. Both sides count, at their
recorded post-game ratings.

Games sort by their end time and alternate between \FitGames{} fitting
games and \HeldGames{} held-out games. The archive contains \EmptyGames{}
games with no counted move. They remain in the game ordering, and those in
the held-out half remain bootstrap units. This preserves the predecessor's
halves. They are held out from this fit, but the predecessor explored the
same archive when it chose the question.

Before any candidate effect was computed, the guard was expanded from the
frozen accuracy cases to every held-out game. The frozen protocol has
content digest \ProtocolDigest; its committed manifest binds that text
and the measured sources. The margin, candidates and other decision rules
were unchanged. The new routine reviews cover all \GuardGames{} held-out games,
including games with no counted move. They supply \PunishEvents{} non-book
reply events in \PunishGames{} games; \PunishVerified{} replies have
verified or terminal evidence.

\section{Candidates and estimation}
Every band comes from the app's conversion and every probability from its
Node engine's \texttt{humanPolicy}. The baseline adds \Base{} for both
people. Candidate A jointly fits separate player and opponent offsets from
\OffsetMin{} to \OffsetMax{} in steps of \OffsetStep{}: \ACount{}
combinations. Each person's offset follows that person into either input
of Maia when the side to move changes.

Candidate B uses the app's rounding and clamping on
$a+b(r-\Anchor)$, with the same intercept grid and slopes \Slopes:
\BCount{} combinations. A line is the smallest family with a slope that
could follow the predecessor's rating thirds. Both slope directions are
allowed. Each family contains the baseline. No grid is expanded after the
run.

The fits maximize mean natural log probability per fitting move. A tie in
A first minimizes squared distance from the baseline pair, then the gap
between its offsets, then the player and opponent offsets. A tie in B first
minimizes absolute slope, then the intercept's distance from the baseline,
then the intercept and slope. Only exact equality is a tie.

The held-out gain is candidate minus baseline mean log probability per
move. Fits stay fixed in \Draws{} shared bootstrap draws of held-out games,
sampled with replacement by the predecessor's generator at seed \Seed.
Each draw sums paired log-probability differences and divides by its move
count. The zero-based percentile index is \QuantileFormula{} on the sorted draws.
An empty denominator makes a comparison inconclusive.

\section{The fixed decision}
A candidate needs a strictly positive fitting gain and a corrected
held-out win. The doubled lower-tail bootstrap p-value for positive gain is
\TailFormula, where $B$ is the draw count. Its finite resolution is
\PResolution; the smallest adjusted value at the first Holm step is
\HolmMinimum. This is a bootstrap approximation, not an exact test.

Holm covers both candidate families at \Alpha. An exact p-value tie puts A
first. A win requires adjusted p strictly below that level and the
\OrdinaryCI\% interval entirely above zero. It also requires a positive
lower end of the step interval: \AdjustedCI\% at the first step
(\StageQuantiles), then \OrdinaryCI\% at the second
(\OrdinaryQuantiles). The second step cannot win if the first fails.

A likelihood winner also needs noninferiority of punish odds. The app's
routine review defines an error as losing at least \ErrorLoss{} points of
winning chances. Its next reply contributes an event if it is non-book and
has an engine best move. The outcome is whether that best reply was played.
The predictor is its Maia probability at the replier's converted band,
with the other input at that same band. This preserves the app's current
punish conditioning.

The measure is the mean Brier score, \BrierFormula. Its difference is
candidate minus baseline; lower is better. Verified replies and all
non-book replies are checked separately in each error direction. Every
one of the \GuardCells{} cells must have its \AdjustedCI\% interval upper
end at most \GuardMargin. These use the same held-game draws, including
zero-event games, and divide by each draw's event count. The correction
covers both candidate guards; all cells are required within a candidate.
Missing evidence or any empty denominator cannot pass.

A sole qualifying candidate is selected. If both qualify, the greater
observed held-out gain is selected.
An exact tie selects B because it needs no owner identity. This comparison
does not establish that the selected candidate significantly beats the
other. If neither qualifies, the baseline stays. The frozen protocol specifies the full rule.

\section{Results}
Candidate A fits \PlayerOffset{} for the player and \OpponentOffset{} for
his opponents. Candidate B fits intercept \Intercept{} and slope \Slope.
Their held-out gains per move are \GainA{} (\OrdinaryCI\% interval
\IntervalA) and \GainB{} (\IntervalB).

\begin{center}\small\FitTable\end{center}
\begin{center}\small\GainTable\end{center}
\begin{center}\small\HolmTable\end{center}

\OutcomeDetail{} \SelectionDetail{}
\Verdict{}

All decisions use full-precision values. The displayed values are rounded.

The following tables retain every guard cell. The ordinary intervals are
\OrdinaryCI\%; the intervals used by the guard are \AdjustedCI\%. The
score table shows the Brier levels as well as their difference.

\begin{center}\small\GuardScoreTable\end{center}
\begin{center}\small\GuardIntervalTable\end{center}
\GuardFinding{}

\section{Beside the verdict}
The following views were specified before the fit and do not decide
adoption. Side and rating groups use the same fitted candidates and held
draws. Rating cuts come from all counted movers, at \RatingCutLow{} and
\RatingCutHigh. The full rating range is \RatingMin{} to \RatingMax.

\begin{center}\small\GroupTable\end{center}
\ExploratoryComparison{}
\SideFinding{}
\RatingFinding{}

First-choice shares are percentages. A tie between top moves is broken by
their order in the app's legal-move policy, as in the predecessor.

\begin{center}\small\TopTable\end{center}

The appendix gives the predecessor's fixed probability bins for each guard
cell and conversion. Their \OrdinaryCI\% Wilson intervals describe observed frequencies.
These bin intervals do not account for clustering within games.
The old bin rule marks a discrepancy only with at least \BinMinimum{}
events and a mean prediction more than \BinAllowance{} beyond the interval.
It is descriptive here and adds no adoption condition.
\BinFinding{}

\section{Measured work}
New work used at most \MaxWorkers{} heavy workers together, each
with a nice increment of \Nice{} from the ordinary-priority coordinator.
The reused reviews from the phase-one timing slice had also run under a
lowered-priority coordinator.
The elapsed span was \ElapsedHours{} hours; recorded session wall
time was \SessionHours{} hours, with \GapHours{} hours of gaps or unmeasured
time. \RuntimeCaveat{}
Session and stage wall times may include scheduling delays or machine
suspension. They do not measure active CPU time.
\RecoveredTimingNote{}

\begin{center}\small\TimingTable\end{center}

The move grid required \MoveCalls{} policy calls and the punish grid
\PunishCalls. Policy files occupy \PolicyMB{} MB. The routine review stage
reused \ReusedReviews{} completed games and newly completed \NewReviews;
the reviewed games contain \ReviewPlies{} plies.
\TimeoutRecoveryNote{}

\begin{center}\small\ReviewTimingTable\end{center}

Completed-review worker time retains the latest completed timing for each
game. Failed timeout time sums the failed child durations in the verified
attempts. A valid review can be written before a timeout, so its recovered
timing can overlap the failed interval. Worker durations also overlap
across concurrent workers. They must not be added to infer stage wall time.
None measures active CPU time.
Stage times count concurrent workers once. The following policy elapsed times instead
sum worker elapsed time on committed rows with durable timing. They omit interrupted rows, serialization
and append overhead; first calls include restarts.
\PolicyTimingCoverageNote{}

\begin{center}\small\PolicyTimingTable\end{center}

The earlier inference projection was \ProjectedHours{} hours, with a
planning allowance of \PlanningHours. It excluded preparation and fitting.
It was a projection from a small slice, not a sustained thermal benchmark.
The actual stage table also includes extraction and fitting.
\ProjectionPriorityNote{}

Reviews use the app's routine service, with a fresh process per game, no
tablebase requests, no eval cache and no review-store writes. This storage
and network configuration differs from the predecessor's CLI. The
instrument remains the app's review pipeline. No app port was written.

\section{Limits and the app}
\PortConclusion{}

The call-site audit records the person each conversion models. A port must
preserve equal-band conditioning for punish and findability, refresh cached
learner levels, and distinguish the bot's explicit Maia band from its model
of the human. Punish and defense currently ignore the profile offset and
need their own plumbing if a conversion is ported.

These are the games of one player and his opponents. Post-game ratings,
the lowest band's merging of weak players, and model mismatch all affect
the fitted conversion. It is not a general conversion between chess sites.
The game bootstrap does not model serial dependence or repeated opponents.
A fresh archive is needed for confirmation independent of the exploration
that motivated this study.

The guard now covers every held-out game, but the verification pass still
selects partly on reply quality. Requiring the all-reply guard helps expose
that selection without removing it. Brier measures calibration and
resolution together. Passing it does not prove calibration in every
subgroup or validate synthetic bot strength.

\appendix
\section{Descriptive punish calibration}
\small
\CalibrationTables
\end{document}
